% sqlite-grdb-internals.tex — SQLite for mobile engineers: pages + B-trees (table B+tree by
% rowid, index B-tree), rollback journal vs WAL (end marks, -wal/-shm, checkpoint, starvation),
% DEFERRED/IMMEDIATE transactions, SQLITE_BUSY, one writer → GRDB DatabaseQueue vs DatabasePool,
% EXPLAIN QUERY PLAN + covering indexes + N+1, FTS5, migrations (user_version, DatabaseMigrator),
% Data Protection classes + WAL sidecars, 0xdead10cc, backups (backup API, VACUUM INTO),
% ValueObservation, Core Data / SwiftData / Room on SQLite.
% Not repeated (neighbours): Core Data contexts/merge/faulting (ios-swift/persistence.tex),
% Core Data model versioning + SwiftData migrations (ios-platform/coredata-migrations-swiftdata.tex).
% Sources (checked 2026-09-26): sqlite.org/wal.html (end mark, 1000-page auto-checkpoint ~4 MB,
% wal-index in shared memory, same host only, checkpoint starvation, persistent journal_mode,
% synchronous=NORMAL not durable on power loss); sqlite.org/lang_transaction.html (DEFERRED /
% IMMEDIATE / EXCLUSIVE, read→write upgrade fails with SQLITE_BUSY, one writer); sqlite.org/
% fileformat2.html (512–65536 pages, table B-tree data only in leaves keyed by 64-bit rowid,
% index B-tree keys only, user_version at offset 60); sqlite.org/pragma.html (page_size 4096
% since 3.12.0, foreign_keys OFF by default, wal_checkpoint PASSIVE/FULL/RESTART/TRUNCATE,
% synchronous=NORMAL for WAL); sqlite.org/eqp.html (SCAN, SEARCH … USING [COVERING] INDEX, USE
% TEMP B-TREE FOR ORDER BY); sqlite.org/lang_vacuum.html (VACUUM INTO = consistent backup of a
% live db); GRDB docs on github.com/groue/GRDB.swift (Concurrency, DatabaseSharing: 0xdead10cc,
% observesSuspensionNotifications, busyMode .timeout; Migrations; ValueObservation; FTS5 needs
% SQLITE_ENABLE_FTS5); Apple Platform Security "Data Protection classes" (Class C default,
% Class A key discarded ~10 s after lock); Apple QA1809 (Core Data SQLite stores default to WAL
% since iOS 7; copy the -wal too).
% @labels: area=ios-platform kind=concept platform=apple topic=data,concurrency,performance
% @tags: sqlite, wal, checkpoint, b-tree, sqlite-busy, grdb, databasepool, explain-query-plan, covering-index, fts5, data-protection, valueobservation
\documentclass[8pt]{extarticle}
\usepackage{printup-sheet}
\usepackage{array}

\lstdefinelanguage{SwiftSheet}{
  morekeywords={protocol,class,final,struct,enum,func,var,let,static,init,for,in,
    if,else,return,guard,self,nil,try,await,async,throws,true,false,
    SELECT,FROM,WHERE,ORDER,BY,CREATE,INDEX,ON,EXPLAIN,QUERY,PLAN},
  sensitive=true, morecomment=[l]{//}, morestring=[b]",
  literate={->}{{\hbox{-}\hbox{>}}}2 {==}{{\hbox{=}\hbox{=}}}2 {??}{{\hbox{?}\hbox{?}}}2
           {!=}{{\hbox{!}\hbox{=}}}2 {--}{{\hbox{-}\hbox{-}}}2}
\lstset{basicstyle=\ttfamily\scriptsize, aboveskip=2pt, belowskip=2pt}

\newcommand\ct[1]{\texttt{#1}}
\newcommand\hd[1]{\par\vspace{3pt}\noindent{\bfseries\color{sheetBlue}#1}\par\vspace{1pt}}

\tikzset{
  pg/.style={draw=sheetGrey, font=\ttfamily\tiny, minimum height=4.2mm, minimum width=6.5mm, inner sep=0.5pt},
  fr/.style={draw=sheetOrange, fill=sheetOrange!12, font=\ttfamily\tiny, minimum height=4.2mm, minimum width=6.8mm, inner sep=0.5pt},
  cmt/.style={draw=sheetGreen!60!black, fill=sheetGreen!25, font=\ttfamily\tiny\bfseries, minimum height=4.2mm, minimum width=5mm, inner sep=0.5pt},
  bt/.style={draw=sheetBlue, fill=sheetBlue!8, font=\ttfamily\tiny, minimum height=4.2mm, inner sep=1pt},
  lbl/.style={font=\tiny, text=black!75, inner sep=1pt, align=center},
  ttl/.style={font=\bfseries\scriptsize, anchor=west},
}

\begin{document}

\sheettitle{SQLite internals for app engineers (+ GRDB)}{ios-platform · folio}

\oneliner{SQLite is a library, not a server: one \textbf{file} of fixed-size \textbf{pages}
(4096 B default) holding \textbf{B-trees} — one per table (rows by \ct{rowid}) and one per index.
Many connections may \textbf{read}, only \textbf{one may write}. In \textbf{WAL} mode commits are
\emph{appended} to \ct{-wal}, readers keep a snapshot (\emph{end mark}), and a \textbf{checkpoint}
copies pages back. Core Data, SwiftData, GRDB and Android Room all sit on it.}

\vspace{2pt}
\noindent\begin{tikzpicture}[sheet]
  % ── B-trees ──
  \node[ttl] at (-0.2,2.55) {Table B+tree (data in leaves) + index B-tree};
  \node[bt, minimum width=16mm] (root) at (2.3,1.95) {root: 400 | 800};
  \node[bt, minimum width=13mm] (l1) at (0.7,1.05) {rows 1–399};
  \node[bt, minimum width=13mm] (l2) at (2.3,1.05) {400–799};
  \node[bt, minimum width=13mm] (l3) at (3.9,1.05) {800–…};
  \foreach \l in {l1,l2,l3} \draw[flow] (root.south) -- (\l.north);
  \node[lbl, anchor=east] at (-0.25,1.05) {leaves:\\full rows};
  \node[bt, draw=sheetGreen!60!black, fill=sheetGreen!10, minimum width=24mm] (ix) at (1.2,0.2) {idx (folderId, rowid)};
  \draw[hot] (ix.east) -- (l2.south);
  \node[lbl, anchor=west, text=sheetOrange] at (2.75,0.35) {rowid $\to$ 2nd lookup};
  \node[lbl, anchor=north west, align=left] at (-0.2,-0.12) {one page read per level: $\approx$3–4 reads for\\millions of rows; \textbf{covering} index = no 2nd lookup};
  \draw[sheetGrey!40] (5.35,2.7) -- (5.35,-0.9);
  % ── WAL ──
  \node[ttl] at (5.5,2.55) {WAL mode — append, snapshot, checkpoint};
  \node[lbl, anchor=east] at (6.8,2.05) {\ct{app.sqlite}};
  \foreach \i/\t in {0/P1,1/P2,2/P3,3/P4,4/P5,5/P6}
    \node[pg, fill=black!4] at (7.2+\i*0.68,2.05) {\t};
  \node[lbl, anchor=east] at (6.8,1.2) {\ct{app.sqlite-wal}};
  \foreach \i/\t/\s in {0/P3/fr,1/P5/fr,2/C1/cmt,3/P3/fr,4/P6/fr,5/C2/cmt,6/P2/fr}
    \node[\s] (w\i) at (7.2+\i*0.68,1.2) {\t};
  \node[fr, draw=sheetRed, dashed, fill=white] at (7.2+7*0.68,1.2) {P4};
  \node[lbl, anchor=west, text=sheetRed] at (12.3,1.2) {writer, uncommitted};
  \node[lbl, anchor=west] at (5.5,0.55) {\ct{-shm}: wal-index};
  \node[lbl, anchor=west, text=sheetGrey] at (5.5,0.3) {page $\to$ newest frame};
  % end marks
  \draw[very thick, sheetBlue] ([xshift=0.34cm]w2.north) -- ++(0,-0.62) node[lbl, below, text=sheetBlue]{reader A};
  \draw[very thick, sheetBlue] ([xshift=0.34cm]w5.north) -- ++(0,-0.62) node[lbl, below, text=sheetBlue]{reader B};
  \node[lbl, anchor=north west, align=left] at (7.3,0.25) {A started before C2: sees P3 of C1, \emph{not} the newer P3;\\B sees C2. Readers never block the writer, and vice versa.};
  % checkpoint
  \draw[->, thick, sheetGreen!60!black] ([xshift=-1mm]w0.north) to[bend left=25] (8.45,1.83);
  \draw[->, thick, sheetGreen!60!black] ([xshift=1mm]w1.north) to[bend right=25] (9.9,1.83);
  \node[lbl, anchor=west, align=left, text=sheetGreen!45!black] at (12.3,2.25) {checkpoint: frames $\to$ db file, only\\up to the \textbf{oldest} reader's end mark;\\auto at \textbf{1000 pages} ($\approx$4 MB)};
  \node[lbl, anchor=west, align=left, text=sheetRed] at (12.3,0.45) {a reader that never ends $\to$\\no checkpoint completes $\to$\\\ct{-wal} grows without bound};
\end{tikzpicture}

\begin{multicols}{2}
\raggedright\small\setstretch{1.0}

\hd{How it works}
\begin{itemize}
  \item \textbf{Pages}: 512–65\,536 B, 4096 default (since 3.12.0). A table is a \textbf{B+tree} keyed by
        the 64-bit \ct{rowid} (\ct{INTEGER PRIMARY KEY} aliases it) — data only in leaves. An
        index is a B-tree of \emph{(columns, rowid)}; \ct{WITHOUT ROWID} tables are index B-trees.
  \item \textbf{Rollback journal} (\ct{DELETE}, the default): original pages copied to \ct{-journal},
        the db written \emph{in place}; commit = delete the journal. A writer committing locks readers out.
  \item \textbf{WAL} (\ct{PRAGMA journal\_mode=WAL}, \textbf{persistent}): commit = append a commit
        frame to \ct{-wal}. Readers use \ct{-shm} (shared memory) — so \textbf{same host only}, no network
        file systems. Big win for apps: UI reads never wait for a background import.
  \item \textbf{Transactions}: \ct{BEGIN} = \textbf{DEFERRED} — no lock until first access; a read that
        later writes must \emph{upgrade}, and fails with \textbf{\ct{SQLITE\_BUSY}} if another connection
        wrote meanwhile. \ct{BEGIN IMMEDIATE} takes the write lock up front (= \ct{EXCLUSIVE} in WAL).
        A read transaction is a \textbf{snapshot}: it never sees later commits.
  \item \ct{SQLITE\_BUSY} = another \emph{connection} holds the lock (busy timeout = wait and retry).
        No timeout set $\to$ fails at once. The real fix: \textbf{one writer} in the whole app.
  \item \textbf{Durability}: WAL + \ct{synchronous=NORMAL} (recommended) syncs only at checkpoint — a
        commit may be lost on \emph{power loss}, never corrupted. \ct{FULL} syncs every commit.
  \item Off by default: \ct{PRAGMA foreign\_keys} (per connection). GRDB turns it on.
\end{itemize}

\hd{GRDB — one writer, many readers}
\begin{itemize}
  \item \textbf{One} \ct{DatabaseQueue} \emph{or} \ct{DatabasePool} per file, for the app's lifetime.
        \textbf{Queue} = one connection, everything serialised. \textbf{Pool} = WAL, a serial
        \textbf{writer} + concurrent \textbf{readers}, each read isolated.
  \item Every \ct{write \{\}} and \ct{read \{\}} is a transaction; writes are serialised (so no
        \ct{SQLITE\_BUSY} inside one process); no writes inside \ct{read}; non-reentrant.
        Escape hatches are named \ct{…WithoutTransaction} / \ct{unsafe…}.
  \item \ct{ValueObservation.tracking \{ db in … \}}: records the \emph{region} it read; any commit
        touching it $\to$ re-fetch. Only committed changes; may \textbf{coalesce}; \ct{removeDuplicates()};
        \ct{.values(in:)} async sequence or \ct{start(in:onError:onChange:)}.
  \item \ct{DatabaseMigrator}: named migrations, each in its own transaction, applied list stored
        in the db; FK checks deferred to commit. \textbf{Never edit a shipped migration}.
        \ct{eraseDatabaseOnSchemaChange} is for development only.
\end{itemize}

\hd{Example — pool, migration, write, observe}
\begin{lstlisting}[language=SwiftSheet]
var cfg = Configuration()
cfg.busyMode = .timeout(5)        // app-group peers: wait, don't fail
let db = try DatabasePool(path: url.path, configuration: cfg)
var m = DatabaseMigrator()
m.registerMigration("v1") { db in
  try db.create(table: "note") { t in
    t.autoIncrementedPrimaryKey("id")
    t.column("folderId", .integer).notNull()
    t.column("title", .text).notNull()
    t.column("updatedAt", .datetime).notNull() }
  try db.execute(sql: """
    CREATE INDEX idx_note_folder ON note(folderId, updatedAt, title)""") }
try m.migrate(db)                  // each migration = one transaction
try await db.write { db in try Note(folderId: 1, title: "Hi").insert(db) }
let obs = ValueObservation.tracking { db in    // region: note table
  try Note.filter(Column("folderId") == 1).fetchAll(db) }
for try await notes in obs.values(in: db) { render(notes) }
\end{lstlisting}
\begin{itemize}
  \item Tests: \ct{DatabaseQueue()} with no path = in-memory db — run the \emph{real} migrator on it,
        plus golden db files written by old app versions.
\end{itemize}

\columnbreak

\hd{Query plans, indexes, N+1}
\begin{lstlisting}[language=SwiftSheet]
EXPLAIN QUERY PLAN SELECT title FROM note
  WHERE folderId = ? ORDER BY updatedAt;
-- no index:    SCAN note + USE TEMP B-TREE FOR ORDER BY
-- (folderId):  SEARCH note USING INDEX ... (folderId=?)
--              + USE TEMP B-TREE FOR ORDER BY
-- (folderId, updatedAt, title):
--   SEARCH note USING COVERING INDEX idx_note_folder (folderId=?)
\end{lstlisting}
\begin{itemize}
  \item \textbf{SCAN} = every row; \textbf{SEARCH} = B-tree seek; \textbf{COVERING} = no table lookup.
        Order: \textbf{equality} columns first, then the \ct{ORDER BY}/range one. \ct{TEMP B-TREE} = a
        sort you could index away. Every index costs each write.
  \item \textbf{N+1}: 1 query for folders + 1 per folder. Fix: \ct{JOIN}, \ct{IN (…)}, or GRDB
        \ct{including(all:)} (one query per association, not per row).
  \item \textbf{FTS5}: \ct{USING fts5(title, body)} + \ct{MATCH} + \ct{ORDER BY rank} (bm25); GRDB
        \ct{FTS5()} + \ct{synchronize(withTable:)}. Needs \ct{SQLITE\_ENABLE\_FTS5} (GRDB docs:
        custom build). \ct{LIKE '\%x\%'} is a scan.
\end{itemize}

\hd{On iOS — files, locks, backups}
\begin{itemize}
  \item \textbf{Three files}: db + \ct{-wal} + \ct{-shm}. Copy, move, back up or protect only the db
        $\to$ lost commits. Put them in \textbf{one directory} and treat it as a unit.
  \item \textbf{Data Protection}: default Class C (\ct{…UntilFirstUserAuthentication}). \ct{.complete}
        (A): key gone $\approx$10\,s after lock $\to$ background writes fail (\emph{disk I/O error}).
  \item \textbf{\ct{0xdead10cc}}: suspended holding a file lock in a \textbf{shared container} (app
        group) $\to$ killed. GRDB: \ct{observesSuspensionNotifications} + \ct{Database.suspendNotification}.
  \item \textbf{Backup} live: backup API (GRDB \ct{backup(to:)}) or \ct{VACUUM INTO 'copy.sqlite'}
        (consistent, compacted). Never \ct{cp} mid-transaction.
  \item \textbf{Migrations} without a library: \ct{PRAGMA user\_version} (int at header offset 60) —
        read, run steps above it in one transaction, write it back.
  \item \textbf{Core Data} SQLite stores default to WAL since iOS 7; \textbf{SwiftData} is Core Data
        underneath. \textbf{Android}: Room = SQLite + compile-time-checked DAOs, \ct{Flow} observation.
\end{itemize}

\hd{Traps}
\begin{itemize}
  \trap{Two pools on one file — two writers, \ct{SQLITE\_BUSY}, observations miss commits.}
  \trap{Long read transaction (a paused cursor, a huge export) $\to$ checkpoint starvation, \ct{-wal}
        of gigabytes, slower reads.}
  \trap{Read-then-write in a DEFERRED transaction $\to$ \ct{SQLITE\_BUSY} at the upgrade; write
        transactions from the start (GRDB \ct{write}).}
  \trap{``Add an index'' without \ct{EXPLAIN QUERY PLAN} before \emph{and} after.}
  \trap{Deleting or overwriting the db but not its \ct{-wal}, or unlinking a file still open —
        a mispaired log is a documented corruption path. So is linking \textbf{two SQLite copies}
        (system + a bundled one): one's \ct{close()} can drop the other's POSIX locks.}
\end{itemize}

\hd{Remember}
\textbf{Pages $\to$ B-trees $\to$ one writer.} WAL: \emph{append, snapshot, checkpoint} —
and the oldest reader decides when the log can shrink.


\end{multicols}

\noindent{\footnotesize\color{sheetGrey}\textit{Related:} persistence (store choice, Core Data
contexts) · coredata-migrations-swiftdata · data-access-patterns · keychain-secure-enclave (Data
Protection) · app-extensions (app groups) · background-execution}

\end{document}
